\documentclass[11pt]{article}

\usepackage{cmap}
\usepackage[T2A]{fontenc}
\usepackage[utf8]{inputenc}
\usepackage[english,russian]{babel}

\usepackage{amsmath, amssymb, amsthm, mathtools}
\usepackage[margin=2.2cm]{geometry}
\usepackage{enumitem}
\usepackage{microtype}
\usepackage{xcolor}
\usepackage{titlesec}

\titleformat{\section}{\large\bfseries}{\thesection.}{0.6em}{}
\titlespacing*{\section}{0pt}{10pt}{4pt}

\setlength{\parindent}{0pt}
\setlength{\parskip}{2pt}

\tolerance=2000
\emergencystretch=2em

\theoremstyle{definition}
\newtheorem{problem}{Задача}

\newtheoremstyle{starred}{}{}{}{}{\bfseries}{.}{.5em}%
  {\thmname{#1}\thmnumber{ #2}$^{*}$\thmnote{ (#3)}}
\theoremstyle{starred}
\newtheorem{starproblem}[problem]{Задача}
\theoremstyle{definition}

\theoremstyle{remark}
\newtheorem*{remark}{Замечание}
\theoremstyle{plain}
\newtheorem*{theorem*}{Теорема}
\theoremstyle{definition}

\DeclareMathOperator{\tr}{tr}
\DeclareMathOperator{\rank}{rank}
\DeclareMathOperator{\conv}{conv}
\DeclareMathOperator{\aff}{aff}
\DeclareMathOperator{\diag}{diag}
\DeclareMathOperator{\dom}{dom}
\DeclareMathOperator*{\argmin}{argmin}
\DeclareMathOperator{\sign}{sign}
\newcommand{\D}{\mathrm{D}}
\newcommand{\R}{\mathbb{R}}
\newcommand{\Sym}{\mathbb{S}}
\newcommand{\norm}[1]{\left\lVert #1 \right\rVert}
\newcommand{\iprod}[2]{\langle #1,\, #2 \rangle}
\newcommand{\ones}{\mathbf{1}}
\newcommand{\pd}{\partial}

\begin{document}

\begin{center}
{\Large\bfseries Домашнее задание \#2}\\[3pt]
{\normalsize Методы оптимизации, ФПМИ, осенний семестр 2026--2027}\\[6pt]
{\normalsize\itshape Дедлайн: 8.11.2026, 23:59}\\[3pt]
{\normalsize\itshape Максимальный балл: 50}
\end{center}

\vspace{0.3cm}

\section{Субдифференциал}

\begin{problem}[Субдифференциал расстояния, 3 балла]
Найдите субдифференциал функции $f(x) = \operatorname{dist}(x, C) = \min_{y \in C}\norm{x - y}_2$
для непустого замкнутого выпуклого $C$ во всех точках. Разберите точки вне $C$, внутри $C$ и на
его границе.
\end{problem}

\begin{problem}[Сумма модулей, 3 балла] 
Пусть $a_1 < a_2 < \dots < a_m$ --- фиксированные числа и $f(x) = \sum_{i=1}^{m} |x - a_i|$,
$f \colon \R \to \R$.
\begin{enumerate}[label=(\alph*), nosep]
  \item Найдите $\pd f(x)$ во всех точках.
  \item Пользуясь теоремой Ферма, найдите все точки минимума $f$. Разберите случаи чётного и
        нечётного $m$.
\end{enumerate}
\end{problem}

\begin{problem}[Выпуклая интерполяция, 4 балла]
Дан набор точек $\{(x_i, f_i, g_i)\}_{i=1}^m$, где $x_i\in\R^n$, $f_i\in\R$, $g_i\in\R^n$. Каждая
тройка обозначает точку, значение функции в ней и значение (суб)градиента. Будем говорить, что для
набора выполняются \emph{условия интерполяции}, если для любых $i, j = 1, \ldots, m$ выполняется
неравенство
\begin{align*}
    f_j \;\ge\; f_i + \iprod{g_i}{x_j - x_i}.
\end{align*}
\begin{enumerate}[label=(\alph*), nosep]
  \item Докажите, что условия интерполяции выполнены тогда и только тогда, когда существует
        выпуклая функция $f \colon \R^{n} \to \R$ с $f(x_i) = f_i$ и $g_i \in \pd f(x_i)$ для всех
        $i$. \emph{Указание:} рассмотрите $\hat f(x) = \max_i\bigl(f_i + \iprod{g_i}{x - x_i}\bigr)$.
  \item Докажите, что $\hat f$ --- наименьшая из таких функций. Приведите пример набора, для
        которого интерполирующая функция не единственна, и пример, когда условия интерполяции
        нарушены, хотя $f_i$ совпадают со значениями выпуклой функции в точках $x_i$.
\end{enumerate}
\end{problem}

\section{Двойственность}

\begin{problem}[Бисопряжённая функция, 5 балла]
Невыпуклую задачу $\min f(x)$ часто заменяют выпуклой, подставляя вместо $f$ её \emph{выпуклую оболочку} --- наибольшую выпуклую функцию, не превосходящую $f$. Такая замена даёт лучшую возможную выпуклую нижнюю оценку: её минимум не больше $\min f$, а зазор между ними показывает, насколько задача <<невыпукла>>. 

Рассмотрим функцию $f \colon \R^{n} \to \R\cup\{+\infty\}$, у которой есть хотя бы одна аффинная
миноранта: существует аффинная функция $h(x)$ такая, что $h(x)\leq f(x)$ для всех $x\in \R^n$ (выпуклость $f$ не предполагается).
\begin{enumerate}[label=(\alph*), nosep]
  \item Пусть функция $g(x)$ поточечно не превосходит $f(x)$, т.\,е. $g(x) \le f(x)$ для всех
        $x \in \R^{n}$. Докажите, что тогда поточечно выполняется $g^{*}(y) \ge f^{*}(y)$.
  \item Докажите, что $f^{**}$ --- наибольшая выпуклая замкнутая функция, не превосходящая $f$.
        В частности, $f^{**} = f$ для выпуклой замкнутой $f$.
  \item Пусть $f(x) = (x^{2} - 1)^{2}$, $x \in \R$. Найдите $f^{**}(x)$.
  \item У многих важных невыпуклых функций нет информативной глобальной выпуклой миноранты. Пример
        --- число ненулевых компонент (так называемая $\ell_0$-<<норма>>)
        $g(x) = \norm{x}_0 = \sum_{i=1}^{n} [x_i \neq 0]$, где $[\,\cdot\,]$ равно $1$, если
        условие выполнено, и $0$ иначе. Однако содержательную выпуклую аппроксимацию можно
        получить, ограничив функцию на компакт:
        $\tilde g(x) = \norm{x}_0 + \delta_{\{\norm{x}_\infty \le 1\}}(x)$.
        Найдите $g^{*}$, $g^{**}$, $\tilde g^{*}$ и $\tilde g^{**}$ и сделайте выводы.
\end{enumerate}
\end{problem}

\begin{problem}[Субдифференциал спектральной функции, 5 баллов]
Сопряжённая функция $f^{*}(y) = \sup_x\bigl(\iprod{y}{x} - f(x)\bigr)$ описывает $f$ через опорные
гиперплоскости к её надграфику, а субградиент в точке --- это наклон одной из них. Значит, между
$\pd f$ и $\pd f^{*}$ должно быть точное соответствие; оно и служит мостом между негладким анализом
и двойственностью.

Пусть $F \colon \R^{n} \to \R$ выпукла и симметрична (не меняется при перестановке аргументов), а
$f(X) = F(\lambda(X))$ на $\Sym^{n}$, где $\lambda(X)$ --- вектор собственных значений по убыванию.
Выпуклость $f$ доказана в ДЗ \#1. Пользуйтесь без доказательства следующим утверждением.

\begin{theorem*}[Неравенство фон Неймана для собственных чисел]
Для любых $X, Y \in \Sym^{n}$ с векторами собственных значений $\lambda(X), \lambda(Y)\in \R^n$,
упорядоченными по убыванию, $\tr(XY) \;\le\; \iprod{\lambda(X)}{\lambda(Y)} . $ Равенство достигается тогда и только тогда, когда существует ортогональная матрица $Q$ такая, что
$X = Q\diag\bigl(\lambda(X)\bigr)Q^{\top}, \qquad Y = Q\diag\bigl(\lambda(Y)\bigr)Q^{\top}.$

\end{theorem*}

\begin{enumerate}[label=(\alph*), nosep]
  \item Пусть $h$ --- произвольная замкнутая выпуклая функция. Пользуясь утверждением (b) из задачи 4, покажите
        равносильность трёх условий: $y \in \pd h(x)$; \ $h(x) + h^{*}(y) = \iprod{x}{y}$; \
        $x \in \pd h^{*}(y)$.
  \item Докажите, что $f^{*}(Y) = F^{*}(\lambda(Y))$ и $F^*$ является симметричной функцией (см. последнюю задачу в ДЗ  \#1).
  \item Пользуясь пунктами (a) и (b), докажите, что
        \[
        \pd f(X) = \bigl\{Q \diag(\mu) Q^{\top} :\ \mu \in \pd F(\lambda(X)),\
        X = Q\diag(\lambda(X))Q^{\top},\ Q^{\top}Q = I\bigr\}.
        \]
  \item Для $X \in \Sym^{n}$ ядерная норма равна $\norm{X}_{*} = \sum_{i=1}^{n}|\lambda_i(X)|$.
        Найдите $\pd\norm{X}_{*}$.
\end{enumerate}
\end{problem}


\begin{problem}[Оболочка Мак-Кормика, 4 балла]
В нелинейных моделях постоянно встречаются произведения переменных: балансовые соотношения вида
<<поток $\times$ концентрация>>, условия смешивания, произведения бинарной и непрерывной
переменных. Каждое такое произведение делает задачу невыпуклой. Стандартный приём --- ввести новую
переменную $z$ для произведения и заменить равенство $z = xy$ выпуклым множеством, содержащим его
график; чтобы релаксация была как можно точнее, берут выпуклую оболочку графика.

Пусть
\[
G = \{(x,y,z) \in \R^{3} : z = xy,\ x, y \in [0,1]\},
\qquad
P = \{(x,y,z) : z \ge 0,\ z \ge x{+}y{-}1,\ z \le x,\ z \le y\}.
\]
\begin{enumerate}[label=(\alph*), nosep]
  \item Покажите, что $G$ невыпукло и что $G \subseteq P$.
  \item Вычислите опорную функцию $h_G(c) = \sup_{(x,y,z) \in G} \iprod{c}{(x,y,z)}$.
  \item Выпишите двойственную задачу к задаче линейного программирования
        $\max\{\iprod{c}{(x,y,z)} : (x,y,z) \in P\}$ и найдите $h_P(c)$.
  \item Выведите отсюда, что $\conv G = P$.
\end{enumerate}
\end{problem}

\section{Условия оптимальности}

\begin{problem}[Неравенство Канторовича, 3 балла]
\begin{enumerate}[label=(\alph*), nosep]
  \item Пусть $a \in \R^{n}_{++}$ ($n\geq 2$) имеет упорядоченные по убыванию компоненты
        $a_n \ge \dots \ge a_1 > 0$, а $b \in \R^{n}_{++}$ определён как $b_k = 1/a_k$. Выведите
        условия ККТ для задачи
        \[
        \min_{x \in \R^{n}} \ -\log(a^{\top}x) - \log(b^{\top}x)
        \qquad\text{при}\qquad x \ge 0,\ \ \ones^{\top} x = 1,
        \]
        и покажите, что $x = \bigl(\tfrac12, 0, \dots, 0, \tfrac12\bigr)^{\top}$ --- её решение.
  \item Пусть $A \in \Sym^{n}_{++}$ с собственными значениями
        $\lambda_n \ge \dots \ge \lambda_1 > 0$. Применив пункт (a) к вектору
        $a = (\lambda_n, \dots, \lambda_1)^{\top}$, докажите неравенство Канторовича
        \[
        2\bigl(u^{\top} A u\bigr)^{1/2}\bigl(u^{\top} A^{-1} u\bigr)^{1/2}
        \;\le\; \sqrt{\frac{\lambda_n}{\lambda_1}} + \sqrt{\frac{\lambda_1}{\lambda_n}},
        \qquad \norm{u}_2 = 1 .
        \]
\end{enumerate}
\end{problem}

\begin{problem}[Проекция на симплекс, 4 балла]
Рассмотрим задачу
$\min\bigl\{\norm{x - y}_2^{2} : \ones^{\top} x = 1,\ x \ge 0\bigr\}$.
\begin{enumerate}[label=(\alph*), nosep]
  \item Выведите условия ККТ, упростите их и сведите задачу к нахождению одной двойственной
        переменной $\nu \in \R$, отвечающей ограничению-равенству.
  \item Постройте алгоритм нахождения оптимального $\nu$ с учётом полученной зависимости $x(\nu)$
        и оцените его сложность.
\end{enumerate}
\end{problem}

\begin{problem}[Подзадача метода доверительной области, 5 баллов] Многие итерационные методы на каждом шаге заменяют функцию квадратичной моделью и минимизируют
её. Чтобы модель не уводила слишком далеко, шаг ограничивают шаром --- областью, где модели ещё
можно доверять. В невыпуклой задаче матрица модели бывает неопределённой, так что подзадача
невыпукла. Тем не менее она решается точно, и всё держится на условиях оптимальности.

Рассмотрим
\[
\min_{x \in \R^{n}} \ q(x) = x^{\top} A x + 2 b^{\top} x
\qquad\text{при}\qquad \norm{x}_2 \le \Delta,
\]
где $A \in \Sym^{n}$ произвольна (не обязательно неотрицательно определена), $b \in \R^{n}$,
$\Delta > 0$.
\begin{enumerate}[label=(\alph*), nosep]
  \item Выпишите условия ККТ для этой задачи.
  \item Пусть $x^{*}$ и $\lambda \ge 0$ удовлетворяют
        \[
        (A + \lambda I)\,x^{*} = -b,
        \qquad A + \lambda I \succeq 0,
        \qquad \lambda\bigl(\norm{x^{*}}_2 - \Delta\bigr) = 0,
        \qquad \norm{x^{*}}_2 \le \Delta .
        \]
        Докажите, что тогда $x^{*}$ --- глобальный минимум. \emph{Указание:} проверьте тождество
        $q(x) - q(x^{*}) = (x - x^{*})^{\top}(A + \lambda I)(x - x^{*})
        + \lambda\bigl(\norm{x^{*}}_2^{2} - \norm{x}_2^{2}\bigr)$.
  \item Пусть $A \not\succeq 0$. Покажите, что оптимум лежит на границе шара, и что задача сводится
        к решению одного уравнения $\varphi(\lambda) = \Delta$ относительно
        $\lambda \ge -\lambda_{\min}(A)$, где $\varphi(\lambda) = \norm{(A+\lambda I)^{-1}b}_2$.
        Докажите, что $\varphi$ строго убывает на этом луче.
  \end{enumerate}
\end{problem}

\begin{problem}[Эллипсоид максимального объёма, 5 баллов]
Пусть $C = \{x \in \R^{n} : -1 \le a_i^{\top} x \le 1,\ i = 1,\dots,p\}$, причём $\sum_{i=1}^{p} a_i a_i^{\top} \succ 0$, и пусть
$\mathcal{E} = \{x : x^{\top} Q^{-1} x \le 1\}$ при $Q \in \Sym^{n}_{++}$ --- вписанный в $C$
эллипсоид максимального объёма. Требуется доказать, что $\sqrt{n}\,\mathcal{E} \supseteq C$.
\begin{enumerate}[label=(\alph*), nosep]
  \item Покажите, что $\mathcal{E} \subseteq C$ тогда и только тогда, когда
        $a_i^{\top} Q\, a_i \le 1$ для всех $i$. \emph{Указание:} представьте эллипсоид в виде
        $\mathcal{E} = \{Q^{1/2} y : \norm{y}_2 \le 1\}$.
  \item Объём $\mathcal{E}$ пропорционален $(\det Q)^{1/2}$, поэтому $Q$ есть решение задачи
        \[
        \min_{Q \in \Sym^{n}} \ -\log\det Q
        \qquad\text{при}\qquad a_i^{\top} Q\, a_i \le 1,\ i = 1,\dots,p .
        \]
        Пользуясь условиями ККТ для неё, докажите, что $x \in C$ влечёт $x^{\top} Q^{-1} x \le n$.
\end{enumerate}
\end{problem}

\section{Коническая двойственность}

\begin{problem}[Восстановление сигналов, 5 баллов]
Измерению доступен дискретный сигнал $b \in \R^n$, складывающийся из трёх независимых источников. Один источник $x_1  \in \R^n$ --- разреженный (подавляющее число координат нулевые), и его вклад в $b$ образуется действием линейного преобразования $C_1 \in \R^{n \times n}$. Второй источник $x_2  \in \R^n$ --- ограниченный, и на него действует линейное отображение $C_2  \in \R^{n \times n}$. Третий $x_3  \in \R^n$ --- дискретизация гладкого сигнала, и на него действует матрица $C_3  \in \R^{n \times n}$. Измерение подвержено аддитивному гауссовому шуму с единичной матрицей ковариации, так что нужно минимизировать средне-квадратичное отклонение. В общем случае задача метода наименьших квадратов $\min_{x_1, x_2, x_3} \|\sum_{i=1}^3 C_i x_i - b\|$ недоопределена, однако знание структуры позволяет уточнить её путём добавления регуляризаторов --- функций, которые штрафуют за <<нарушение структуры>>. Таким образом к квадратичной невязке добавляют три штрафа: $\ell_1$ отвечает за разреженность, $\ell_\infty$ не
даёт ни одной компоненте вырваться вперёд, а $\norm{Dx}_3$ с разностной матрицей $D$ штрафует негладкость --- жёстче, чем $\ell_2$, но мягче, чем $\ell_\infty$. Здесь произведение $Dx$ задано разницами $x_i - x_{i+1}$. Получается задача
\[
\min_{x_1, x_2, x_3 \in \R^n} \ \norm{C_1 x_1 + C_2 x_2 + C_3 x_3 - b}_2^{2}
\;+\; \alpha \norm{x_1}_1 \;+\; \beta \norm{x_2}_\infty \;+\; \gamma \norm{Dx_3}_3 ,
\]
где $\alpha, \beta, \gamma > 0$ - параметры регуляризации, $C_1, C_2, C_3\in \R^{n\times n}$ - заданные произвольные матрицы, $D \in \R^{(n-1)\times n}$ - разностная матрица ($[Dx]_i=x_{i+1} - x_i$).

Каждое слагаемое выпукло, но сумма четырёх разнородных норм не укладывается ни в одну стандартную
форму. Обычный приём --- перенести каждое слагаемое в ограничение с помощью надграфиковой
переменной: целевая функция становится линейной, а вся информация о слагаемых переходит в конические ограничения.
\begin{enumerate}[label=(\alph*), nosep]
  \item Квадратичное слагаемое потребует \emph{повёрнутого конуса второго порядка}
        $K_{\mathrm{rot}} = \{(u, v, z) \in \R \times \R \times \R^{m} : 2uv \ge \norm{z}_2^{2},\ u \ge 0,\ v \ge 0\}$.
        Докажите, что он выпуклый и является конусом, и найдите $K_{\mathrm{rot}}^{*}$.
  \item Введите переменные $t_0, t_1, t_2, t_3$ и запишите задачу как минимизацию линейной функции
        при ограничениях вида <<вектор принадлежит конусу>>. 
  \item Постройте двойственную задачу.
\end{enumerate}
\end{problem}

\begin{problem}[Полуопределённая релаксация Max-Cut, 4 балла]
Лагранжева двойственность даёт универсальный способ получать нижние оценки для невыпуклых задач.
Более того, по аналогии с задачей~4, построив двойственную задачу к двойственной, можно получить
выпуклую релаксацию исходной задачи. Посмотрим, как это работает на задаче Max-Cut. На семинаре разбиралась задача $\min\{x^{\top} W x : x \in \{-1,1\}^{n}\}$ при $W \in \Sym^{n}$ и
её лагранжева двойственная
\[
\max_{\nu \in \R^{n}} \ -\ones^{\top}\nu
\qquad\text{при}\qquad W + \diag(\nu) \succeq 0 .
\tag{D}
\]
\begin{enumerate}[label=(\alph*), nosep]
  \item Покажите, что двойственная к (D) имеет вид
        $\min\{\tr(WX) : X \succeq 0,\ X_{ii} = 1\}$.
  \item Докажите, что если оптимальная $X$ имеет ранг $1$, то есть $X = xx^{\top}$, то $x$ --- 
        решение исходной задачи.
  \item Как соотносятся оптимальные значения исходной задачи, задачи (D) и задачи из пункта (a)?
  \item Пусть $v$ --- собственный вектор $X$, отвечающий наибольшему собственному значению, и
        $\hat x = \sign(v)$. Докажите, что $\hat x$ допустим в исходной задаче, и объясните, почему
        значение $\hat x^{\top} W \hat x$ вместе с оптимумом релаксации даёт двустороннюю оценку на
        неизвестный оптимум.
\end{enumerate}
\end{problem}

\end{document}
